\section[PHT3D Exercise 7: Artificial recharge and micropollutant fate]{PHT3D Exercise 7: Artificial recharge and micropollutant fate}

\subsection[Background]{Background}
Observed concentrations of reactive solutes in artificially recharged groundwater systems 
can often be highly dynamic and thus puzzling to interpret. For instance, 
concentration variations might be caused by effects, such as temperature variations, 
changes in groundwater flow rates, variations in the injected water quality, or a combination of those. 
One such example where the latter occurs is a study site in Berlin, Germany \citep[]{greskowiak_2006}, 
where aquifers are  recharged by passing wastewater-impacted surface water through 
engineered infiltration ponds, as schematically illustrated in Figure \ref{fig:berlinpond}. 
At this site, the temperature-dependent degradation of sediment-bound 
organic matter (SOM) was an important control for spatiotemporally varying oxygen concentrations. 
However, in addition, highly variable pond infiltration rates caused 
groundwater flow rates to change accordingly and to contribute 
further to the variability of the redox zonation. The model simulations focused 
initially on matching the data collected from in situ temperature loggers, 
and subsequently, the field-observed redox dynamics were replicated. 
In a final step, the degradation behavior of the micropollutant phenazone, 
a pharmaceutical residue, was investigated. 
\par Through the model-based hypothesis testing and data interpretation 
it could be demonstrated that phenazone degraded in this system solely 
under aerobic aquifer conditions that dominated during winter. 
In contrast, phenazone remained persistent under the nitrate-reducing conditions 
that dominated during summer, while oxygen was depleted rapidly (Figure \ref{fig:berlinpond}).
Constraining the groundwater flow and solute transport processes prior 
to evaluating and calibrating the reactive processes associated with oxygen 
consumption allowed for residence time effects to be distinguished from temperature effects.

\par The hypothesized degradation behaviour was based on the results 
of various earlier laboratory experiments that have demonstrated that 
phenazone is biodegradable under oxic conditions \citep{phdzuelke} 
in the presence of the aerobic bacterium $Phenylobakterium$ $immobile$. 
Batch experiments with sterilised and non-sterilised filter material suggest that aerobic microbial
degradation is the most important removal process.
Sorption of phenazone to filter material, sludge, or soil was
found to be negligible. Importantly, it was also found that the degradation
rate of phenazone is temperature-dependent, as its degradation,
like that of organic matter, relies on microbial metabolism.

\par The model simulations in this exercise will employ and discuss the 
model approach that is described in more detail in \cite{greskowiak_2006}. 
The type of redox-sensitive degradation behaviour that is illustrated in 
this exercise was previously also observed for many other micropollutants such as 
pesticides. 

\begin{figure}[tb]
\centering
\includegraphics[width=1.0\textwidth]{berlin_transect}
\caption[Schematic figure of the modelled infiltration pond transect]{Schematic figure of the modelled infiltration pond transect} 
\label{fig:berlinpond}
\end{figure}

\subsection[Governing kinetic reactions]{Governing kinetic reactions}
The degradation rate $r_{SOM}$ of the sediment-bound organic matter (SOM) that is present throughout the aquifer
is simulated by a Michaelis-Menten type formulation:

\begin{equation}
r_{SOM} = f_{T} \bigg[ r_{ox} \frac{C_{ox}}{K_{ox} + C_{ox}} + r_{nit} \frac{C_{nit}}{K_{nit} + C_{nit}} 
\frac{K_{ox,inh}}{K_{ox,inh} + C_{ox}}+ r_m \bigg]
\label{eq:som_deg}
\end{equation}

where $f_{T}$ is an empirical function that describes the temperature dependent degradation 
of SOM. 
Thereby $C_{ox}$ and $C_{nitr}$ are the concentrations of dissolved oxygen and nitrate, 
respectively, $r_{ox}$, $r_{nitr}$, $r_{mn}$ are the maximum rate constants under oxic, nitrate reducing 
and manganese reduction conditions, respectively, $K_{ox}$, $K_{nitr}$ are the 
corresponding half-saturation constants, $K_{ox,inh}$ is the inhibition constant 
for nitrate reduction in the presence of oxygen. 
The degradation rate of phenazone is computed from:

\begin{equation}
{r_{phena} = r_{phena,max} {\frac{C_{ox}}{(k_{phena,ox} + C_{ox})}} C_{phena} f_T }
\label{eq:phenazone_deg}
\end{equation}

where $r_{phena,max}$ is the degradation rate constant,
$k_{phena,ox}$ is the half-saturation concentration, 
and $C_{phena}$ represents the concentration of phenazone.

The dependence of the reaction rate 
on the temperature T ($^{\circ}$C) is considered through 
the empirical function, $f_T$, as proposed by \cite{oconnel_1990} 
and \cite{kirschbaum_1995}:

\begin{equation}
{f_{T} = exp({\alpha + \beta T (1 - 0.5 {\frac{T}{T_{opt}}})  )}}
\label{eq:temp_dep}
\end{equation}

where $\alpha$ and $\beta$ are fitting parameters and $T_{opt}$ is the optimal
temperature for the degradation reaction. 
$T$ is the water temperature, 

In the PHREEQC database 
the reaction rate expression for phenazone 
as used in \cite{greskowiak_2006} is formulated such that the reaction 
will only proceed in the presence of dissolved oxygen.

There are two parameters, i.e., ${\color{blue} parm(1)}$
and ${\color{blue} parm(2)}$ that can be modified 
within ORTI under the ${\color{blue} Rates}$ Tab. 
If additional parameters should be edited, 
this can be achieved by editing the 
rate expression below accordingly. 

{\color{red}
\begin{quote}
\#\#\#\#\#\#\#\#\#\#\#\#\#\#\#\#\#\#\#\#\#\#\#\#\#\#\\
Phenazone \\
\#\#\#\#\#\#\#\#\#\#\#\#\#\#\#\#\#\#\#\#\#\#\#\#\#\#\\
\# from Greskowiak et al. 2006 (EST)
-start \\
 7 mTmp   = 1000 * tot("Tmp") \\
10  Cphena = TOT("Phenazone") \\
20  Cox = TOT("O(0)") \\
26  a = -1.5; \\
27  b = 0.18; \\
28  Topt = 35 \\
30  fT = exp(a+b*mTmp*(1-0.5*mTmp/Topt)) \\
34  rph\_max = parm(1) \\
36  Khalf = parm(2) \\
50  rate = fT* rph\_max*(Cox/(Khalf+Cox)) * Cphena \\
70 moles = rate * TIME \\
200 save moles \\
-end \\
\end{quote}}

The present exercise consists of three consecutive, increasingly complex stages: 
\begin{itemize}
\item First we construct a steady-state flow and reactive transport model that simulates 
the fate of phenazone under constant temperature condition
\item Subsequently we modify this model to study the impact of the filed-observed dynamic 
temperature changes and associated with those, the temporal variations in the redox zonation 
\item Finally we also consider the highly variable infiltration rates and s
tudy how this affects the dynamics of the biogeochemical processes.
\end{itemize}


\subsection[Steady flow model]{Steady flow model}
To set up the steady state flow model
\begin{itemize} 
\item Create a new model an set the model up as a 
{\color{blue}confined} 
{\color{blue}Xsection} model within the
{\color{blue}modflow series}.

\item Set the geometry of the model domain with an extension in the x-direction from 0-125 m 
and in z-direction from 0-41 m. 
The grid discretization can be uniform with dx = 6.25m and dz = 4.1m.

\item To assist with the model construction, 
position a map of the cross section (Background\_image.png) to use it as a background image during the model setup. 
To do this use the {\color{blue}Parameters} $\rightarrow$ {\color{blue}Model} $\rightarrow$ {\color{blue}Map} button

\item The total simulation period considered is one year by going to 
{\color{blue}Parameters} $\rightarrow$ {\color{blue}Model} $\rightarrow$ {\color{blue}time} button
and setting the maximum time to 360 days. Set the timestep length to 5 days, which will cause 
the model to force a reaction simulation every 5 days.

\item Position a boundary condition for modflow ({\color{blue}bas.3}) at the right side (east) 
of the cross section from z=0.25m to 33.25m. 

\item Add at the same location a fixed head ({\color{blue}bas.5}) boundary condition at H = 35 m. 
Set the head for the background at 40 m, which will then be used as the starting head of the model.
\end{itemize}

The pond will be considered as a discrete recharge area, while the recharge 
in other areas can be neglected. 
In order to set this recharge area that represents the pond, 

\begin{itemize}
\item Go into the {\color{blue}Addin menu} $\rightarrow$ {\color{blue}modflow modules} 
and select the {\color{blue}RCH} module. To see the module in the interface just click again on modflow. 

\item In this first phase of the exercise we will apply a steady recharge 
rate of 1.4 m/d under {\color{blue}RCH} $\rightarrow$ {\color{blue}rch.2} 
keep the background value at 0 and draw a horizontal line at the model 
top at the location of the pond, i.e., from x=0m to x = 53m, and provide 
a value of 1.4 (m/d). 

\item As the infiltration rate is very high, 
we need to increase the number of internal time steps, 
in order to achieve an adequate numerical solution of the reactive transport problem. 
Click the {\color{blue}Parameter} Button under {\color{blue}Flow} 
and chose {\color{blue}DIS} $\rightarrow$ {\color{blue}dis.8} 
and change the {\color{blue}NSTP} value to 2. 
This results in 2 time steps within each stress period. 

\item Set the hydraulic conductivity to 50 m/d by going to {\color{blue}LPF} $\rightarrow$ {\color{blue}lpf.8} 
and writing 50 in the background box before clicking {\color{blue}ok}. 
This will apply to the whole model. 

\item The vertical hydraulic conductivity 
is set by defining a ratio for Kh/kv. 
To do this go first to {\color{blue}Flow} $\rightarrow$ 
{\color{blue}Parameters} $\rightarrow$ {\color{blue}LPF} $\rightarrow$ {\color{blue}lpf.5} 
and set the value to a ratio kh/kv. 
Then in {\color{blue}Spatial attributes} $\rightarrow$ {\color{blue}modflow} $\rightarrow$ {\color{blue}lpf.9} 
set the value of the ratio of 1 for the whole model (use the background value)

\item Now run MODFLOW.

\end{itemize}

\subsection[Steady flow, chemistry and constant temperature (Stage 1)]
{Steady flow, chemistry and constant temperature (Stage 1)}

To integrate the reactive transport processes into our model 

\begin{itemize}
\item You will first specify the advection scheme that will be used. 
We will select the MOC scheme in 
{\color{blue}Parameters} $\rightarrow$ {\color{blue}Transport} $\rightarrow$ {\color{blue}Parameters} 
$\rightarrow$ {\color{blue}ADV} $\rightarrow$ {\color{blue}adv.1 major data}.
 
\item Subsequently the dispersivity is defined under
{\color{blue}Spatial attributes} $\rightarrow$ {\color{blue}mt3dms} $\rightarrow$ {\color{blue}dsp.2} 
for the longitudinal dispersion, where we will keep the default value of 1 m.

\item Furthermore, we will also leave the ratio of the vertical to 
the longitudinal dispersivity at 0.05. 

\item The porosity (0.35) is set in the 
{\color{blue}Spatial attributes} $\rightarrow$ {\color{blue}mt3dms} $\rightarrow$ 
{\color{blue}btn.11} $\rightarrow$ {\color{blue}porosity}. 

\item Import the pht3d\_datab.dat via the import button under Chemistry. 
You will get a long list of aquoeus components under 
{\color{blue}(Add\_Chemistry} Button $\rightarrow$ {\color{blue}Solutions}. 
However, we will need to activate only a few of those.

\item Check the file {\color{blue}Berlin\_data\_for\_model.xls} sheet stage 1. 
Copy the species values for the {\color{blue}Backgrd} and the {\color{blue}Solu1} column 
from the spreadsheet into the corresponding columns of the ORTi3D interface. 
Check the values, and click those active where at least one value of 
the {\color{blue}Backgrd} and {\color{blue}Solu1} is non-zero. 

\item Note that for numerical reasons 
the species Tmp carrying the Temperature is given in milli $^\circ$C.
This means that, for example, the input value representing a 
temperature of 19 $^\circ$C would be 0.019.

\item Under {\color{blue}Rates}, all {\color{blue}RATES} are imported 
from the {\color{blue}pht3d\_datab.dat} database. 
Select only {\color{blue}Orgc\_sed} and {\color{blue}Phenazone}. 
Also, select {\color{blue}Orgc\_sed} to act as an immobile (IM) species
to reflect the fact that this organic matter is attached to sediments.
\end{itemize}

The {\color{blue}Orgc\_sed} species is here not 
really meant as the concentration of sedimentary organic matter. 
Rather, it is used in here as a reactivity parameter that defines the 
overall reactivity of the organic matter in the aquifer, and linked to the 
consumption of terminal electron acceptors via {\color{blue}RATES} in the database. 
It was assumed that that there is always a sufficient amount of degradable 
organic matter such that it does not 
become rate-limiting. The parameters {\color{blue}parm1} to {\color{blue}parm5} 
correspond to the parameters that are used in the database. 

\begin{table}[b]
\caption{Reaction parameters in the included reaction rate expressions.}
\begin{center}
\begin{tabular}{cccccc}
\toprule
                 & Parm(1) & Parm(2)                  & Parm(3)    & Parm(4)     & Parm(5)        \\
 SOM / Org\_sed  & $r_{ox}$ & $K_{ox}$ = $K_{ox,inh}$  & $r_{nitr}$ &  $K_{nitr}$ & $r_{mn}$ \\
 Phenazone       & $r_{phena,max}$ & $K_{phena,ox}$        &         &         &          \\
 \bottomrule
\end{tabular}
\label{parms_berlin}
\end{center}
\end{table}

All kinetic reaction and temperature parameters are given 
in Table \ref{parm_values_berlin}. 
Note that the constants $T_{opt}$, $\alpha$ and $\beta$ 
were hard-coded in the {\color{blue}pht3d\_database.dat} database 
and don't need to be specified via the ORTi3D user interface.


\begin{table}[b]
\caption{Parameter values for kinetically controlled reactions}
\begin{center}
\begin{tabular}{ccc}
\toprule
  Parameter             &  Value    &    Unit   \\  $\times$ $10^{-3}$ mol/L$_{bulk}$
$r_{ox}$                & 1.32 $\times$ $10^{-5}$  & mol/L/d   \\
$K_{ox}$ = $K_{ox,inh}$ & 1.00 $\times$ $10^{-5}$  & mol/L     \\
$r_{nitr}$              & 3.31 $\times$ $10^{-6}$  & mol/L/d   \\
$K_{nitr}$              & 1.00 $\times$ $10^{-5}$  & mol/L     \\
$r_{mn}$                & 7.70 $\times$ $10^{-8}$  & mol/L/d   \\
$r_{phena,max}$         & 12        & 1/d       \\
$K_{phena,ox}$          & 4.00 $\times$ $10^{-4}$  & mol/L     \\
$T_{opt}$               & 35        & $^{\circ}$C   \\
${\alpha}$            & -1.5      & -             \\
${\beta}$             & 0.18      & 1/$^{\circ}$C \\
 \bottomrule
\end{tabular}
\label{parm_values_berlin}
\end{center}
\end{table}

\begin{itemize}
\item For the rate {\color{blue}Orgc\_sed} the stoichiometry of the reactive constituents 
associated with the reaction has to be entered under 
formula: Enter {\color{blue}Orgc\_sed -1.0 CH2O 1.0 NH3 0.15}.

\item Subsequently enter the formula {\color{blue}Phenazone -1} to be 
associated with the rate of {\color{blue}Phenazone} degradation.

\item Go to {\color{blue}Phases} and click  {\color{blue}Calcite} and provide a background concentration
under {\color{blue}Backg\_Moles} of 1.7 and for {\color{blue}Pyrolusite} 
add a concentration of 3.0 $\times$ $10^{-3}$ under {\color{blue}Backg\_moles} and similarly enter 0 for 
{\color{blue}Rhodochrosite}. 

\item Now we have to assign the aqueous solutions 
that will act as initial and boundary conditions in the model. 
Choose {\color{blue}Pht3d} $\rightarrow$ {\color{blue}PH} $\rightarrow$ {\color{blue}ph.3 Initial concentrations}. 
Make sure that {\color{blue}Backg.} is set to 0, which means that 
solution 0 is used to define the initial water composition.

\item Now choose {\color{blue}ph.5 Recharge}. 
Define a zone that coincides with the infiltration area 
and enter solution number 1 into the right recharge box. 

\item Finally, everything is ready to run PHT3D.
  
\item To visualise the simulation results of, for example, 
temperature, oxygen, nitrate and phenazone at monitoring wells, 
set observations zones {\color{blue}Spatial attributes} 
$\rightarrow$ {\color{blue}observation} $\rightarrow$ {\color{blue}obs.1} 
create lines that span the screened distance of the wells with well names. 
Then, in the panel {\color{blue}results} $\rightarrow$ {\color{blue}observation}, 
select time series graphs and chemistry. 
In the dialog select the wells and then a layer (it can be any layer) 
and click twice on {\color{blue}Apply}. 

\item To better compare the graphs, they may be exported and then 
the txt file can be opened in excel or further processed by python or matlab.

\end{itemize}

\subsection[Seasonal temperature dynamics (Stage 2)]
{Seasonal temperature dynamics (Stage 2)}

\begin{itemize}
\item First, under {\color{blue}Chemistry} $\rightarrow$ 
{\color{blue}Parameters} $\rightarrow$ {\color{blue}PH} 
$\rightarrow$ {\color{blue}ph.6} $\rightarrow$ 
{\color{blue}specific options} and increase the {\color{blue}nb solutions parameter} to {\color{blue}12}.
This will allow to define 12 individual aqueous solutions (default is max 4).  

\item After that re-import the database file {\color{blue}pht3d\_datab.dat}.

\item Under the {\color{blue}Ad\_chemistry} feature, copy the composition from 
solution 1 ({\color{blue}solu1}) to solutions 2 ({\color{blue}solu2}) - solution 12 ({\color{blue}solu12}). 

\item Now you can enter or copy and paste the varying temperatures ({\color{blue}Tmp}) 
for {\color{blue}solu1} to {\color{blue}solu12}. 
The values are provided in the Excel spreadsheet {\color{blue}Stage 2} (indicated in yellow) 
into the ORTi Interface to the species Tmp. 

\item Now choose {\color{blue}ph.5 Recharge}. 
Choose the zone for infiltration area, make it transient, and copy and paste the solution time series 
(time and solution number) given in Sheet {\color{blue}Stage 2} 
(marked in yellow, without the titles) into the right recharge box.
 
\item Run PHT3D and study how the breakthrough curves at the observation wells have changed.
 
\end{itemize}

\subsection[temperature retardation due to heat exchange with the aquifer matrix]{Temperature retardation due to heat exchange with the aquifer matrix}

As now the temperature varies with time in the infiltration water, 
the generated temperature signal in the aquifer will be retarded due to 
heat exchange with the aquifer matrix. The retardation factor depends on the 
thermal characteristics of the porous medium is approximately R = 2.1 in the Berlin aquifer. 
This can be simulated with the {\color{blue}RCT} reaction package, specified 
under MT3DMS (this will automatically also assigned to PHT3D). 
To do so, we have to provide MT3DMS with the apparent density of the medium 
and the Kd to calculate the retardation factor R (R = 1+$rho_{a}$*Kd/porosity). 
With a porosity of 0.35 and a bulk density of 1722 kg/m3, this gives a Kd of 2.2 $\times$ $10^{-4}$. 

\begin{itemize}
\item To set these values you will have to add the {\color{blue}RCT} module to MT3DMS by
(menu {\color{blue}Addin} {\color{blue}mt3dms modules}).

\item then to go to the dialog 
{\color{blue}Parameters} {\color{blue}mt3dms} $\rightarrow$ {\color{blue}Parameters} 
$\rightarrow$ {\color{blue}RCT} $\rightarrow$  {\color{blue}rct.1} 
and set the sorption to linear. 

\item Once this is done go to 
{\color{blue}Spatial attributes} $\rightarrow$ {\color{blue}Mt3dms} $\rightarrow$ 
{\color{blue}Rct} $\rightarrow$  {\color{blue}rct.2a density} 
and set the value at 1722 for the background (i.e., everywhere). 

\item Then go to {\color{blue}Parameters} $\rightarrow$ {\color{blue}Transport} 
and click on {\color{blue}SPC}, 
select {\color{blue}pht3d} in the choice box and write {\color{blue}Tmp} in the species box. 

\item After that, go to {\color{blue}Parameters} $\rightarrow$ {\color{blue}transport} $\rightarrow$ 
{\color{blue}Rct} and you will get a new dialog box with 4 parameters. 

\item Now enter a Kd value of 2.2 $\times$ $10^{-4}$ in the {\color{blue}SP1} column.

\item Re-run PHT3D 

\end{itemize}

Inspect how this changes the time-concentration curves at the observation wells. 
You will notice that now the temperature breakthrough is slightly shifted towards later times, 
as the retardation of the heat transport also affects the redox system and thus ultimately 
also the phenazone breakthrough behaviour.

\subsection[Seasonal temperature and variable infiltration (Stage 3)]
{Seasonal temperature and variable infiltration dynamics (Stage3)}

In this part of the exercise the recharge at the pond (infiltration rate) will changed to a time-varying input
to reflect the highly dynamic changes in infiltration rates that occurred at the field site.
The 1-year long simulation period is therefore split into 12 monthly stress-periods. 
The employed rates are provided in the Excel spreadsheet {\color{blue}Stage 3}.

\begin{itemize} 
\item Select the two columns time and recharge (marked in yellow, without the titles) 
and copy them into ORTi to {\color{blue}RCH} $\rightarrow$ {\color{blue}rch.2}, 
select the recharge zone and click on {\color{blue}value}, which triggers the {\color{blue}zone dialog} to open. 
In this dialog select the {\color{blue}transient} option. A grid will appear in {\color{blue}values}.
Now paste the recharge values that you copied from the spreadsheet into this grid. 

\item Click the {\color{blue}Parameter} Button under {\color{blue}Flow} 
and select {\color{blue}DIS} $\rightarrow$ {\color{blue}dis.8} and change the simulation mode from 
{\color{blue}steady-state} to {\color{blue}transient} simulation.

\item You now need to re-run MODFLOW. You can then check 
how the heads values and gradients are dynamically changing in the model.

\item Finally re-run PHT3D and inspect how the transient recharge affected the breakthrough behaviour.

\end{itemize}

